%!TEX root=main.tex
\chapter{Introduction} \label{cap:introducao}

Fully Homomorphic Encryption (FHE)~\cite{martins2017survey} is a cryptographic technique that enables computations, such as addition and multiplication, to be performed directly on encrypted data without requiring access to the underlying plaintext. This property enables a user to encrypt its data using an FHE scheme and securely outsource the computation to a third-party or cloud server. The server can then perform the required computations over the encrypted data and return the result in encrypted form. Finally, the user can decrypt the result using its secret key and obtain the corresponding plaintext output.

The main advantage of FHE is that it enables data to remain confidential throughout the computation, thereby reducing the need to trust the computing server. This property is particularly relevant for applications involving sensitive information and can help organizations comply with legal and regulatory requirements concerning the processing and transmission of personal data. For example, consider a facial-recognition system deployed on an edge device that needs to send a person's facial template to a remote server responsible for performing the recognition. In scenarios where transmitting such sensitive biometric information in plaintext is not permitted, FHE or other privacy-preserving techniques can be employed to allow the recognition process to be performed while protecting the underlying data~\cite{martin2025hydia}.

Several FHE schemes have been proposed in the literature, including TFHE~\cite{chillotti2020tfhe}, BFV~\cite{fan2012somewhat}, BGV~\cite{brakerski2014leveled}, and CKKS~\cite{cheon2017homomorphic}. These schemes differ in their underlying constructions and, in particular, in the type of data and computations they are designed to efficiently support. TFHE is primarily designed for bit-level computations, whereas BFV and BGV are commonly used for computations over integers. CKKS, in contrast, is designed to efficiently support approximate arithmetic over real and complex values, making it particularly suitable for applications involving numerical computations, such as machine learning and signal processing.

This work focuses on the CKKS scheme due to its suitability for approximate numerical computation. In particular, its ability to efficiently perform arithmetic operations over real and complex values makes it a promising candidate for privacy-preserving applications involving numerical workloads, including deep-learning inference.

%%%%%%%%%%%%%%%%%%%%%%%%%%
\section{Motivation}

One of the main challenges of CKKS implementations is its high computational overhead on both the client and server sides. Consequently, several works have investigated techniques to reduce this cost through optimizations at the algorithmic, software, and hardware levels.

At the algorithmic level, \cite{cheon2018full} introduced a Residue Number System (RNS) variant of CKKS, enabling more efficient arithmetic over ciphertexts. \cite{halevi2018faster} proposed faster key-switching and linear transformation techniques, which were subsequently further optimized by \cite{bossuat2021efficient} and \cite{yan2026faster}. \cite{cheon2023homomorphic} introduced new ciphertext multiplication algorithms that reduce modulus consumption and, consequently, the computational and memory costs associated with homomorphic operations. Bootstrapping, which remains the most computationally expensive operations in CKKS, has also received considerable attention. Several works have proposed algorithms and optimizations to improve its efficiency, including \cite{cheon2018bootstrapping,chen2019improved,cheon2018faster,kim2022evalround,sung2025evalround+,cheon2025ship,kim2025overmodraise}.

At the software level, several libraries provide implementations of FHE schemes, including CKKS, targeting general-purpose CPUs. Examples include \cite{al2022openfhe,sealcrypto,boemer2021intel,jung2021accelerating,halevi2014algorithms}. These libraries facilitate the integration of software and hardware optimizations into individual CKKS operations, and some already provide support for vectorized instruction sets such as AVX-512. In addition to CPU-based implementations, several works have explored the use of GPUs to accelerate computationally intensive CKKS operations \cite{wang2025chameleon,wang2023he,chen2026hi}.

At the hardware level, dedicated accelerators have been proposed to improve the performance of CKKS operations on both the client and server sides. Client-side accelerators target operations such as encryption/decryption and key generation \cite{yune2025abc,wang2024compact,azad2023rise,krieger2024aloha}, while server-side accelerators focus on the computationally intensive operations involved in homomorphic evaluation \cite{fan2025fast,kim2022ark,kim2023sharp,putra2023strix}. These approaches can provide substantial performance and/or energy-efficiency improvements, but generally rely on specialized hardware, GPUs, or high-end CPUs.

However, none of these works present an highly efficient software implementation of CKKS for edge devices with limited computational resources and without FHE hardware acceleration.

%%%%%%%%%%%%%%%%%%%%%%%%%%
\section{Objectives}

This work aims to develop a highly efficient C-language implementation of the CKKS client-side operations targeting edge devices, with a particular focus on the ARM Cortex-A55~\cite{arm2023cortexa55} from the ARMv8 architecture. ARMv8 based processors are widely used in edge computing platforms and typically provide 32 128-bit Single-Instruction Multiple-Data (SIMD) registers through the ARM NEON extension, as well as dedicated cryptographic instructions for accelerating operations based on the Advanced Encryption Standard (AES). Moreover, the ARM Cortex-A55 features four cores, which can be leveraged to exploit multi-threading and improve the performance of the proposed application. Our implementation aims to exploit these architectural features along with multiple software optimizations to efficiently execute the CKKS scheme.

Furthermore, to provide a complete end-to-end implementation and evaluate the interaction between the two parties, we implement the server-side CKKS operations as well, with the exception of bootstrapping. Although the primary focus of this work is the client side, implementing the server-side operations allows us to provide a software framework for evaluating CKKS applications on edge platforms.

%%%%%%%%%%%%%%%%%%%%%%%%%%
\section{Justification}

Despite the advances at hardware and software level, there remains a gap in providing a complete and highly efficient software implementation of CKKS specifically targeting edge devices with limited computational resources. This scenario is particularly relevant for privacy-preserving image-processing applications, such as traffic cameras and facial-recognition devices~\cite{martin2025hydia}. Such systems often have stringent latency and cost constraints, while energy consumption may be a secondary concern because the devices can be directly connected to a stable power source. Consequently, deploying a dedicated hardware accelerator, a GPU, or a high-end CPU with AVX-512 support may be impractical due to their additional cost and hardware requirements.

In contrast, many edge image-processing devices already incorporate relatively inexpensive multicore CPUs equipped with smaller vector units, typically supporting 128- or 256-bit SIMD instructions. These platforms provide considerable parallelism without requiring specialized hardware or expensive processors. Therefore, there is an opportunity to exploit the computational resources already available on such devices to develop an efficient CKKS implementation tailored to their architectural constraints. This motivates our work, which focuses on optimizing CKKS operations for multicore edge CPUs, with particular emphasis on exploiting SIMD parallelism and the available computational resources while avoiding the hardware and cost overhead associated with GPUs, dedicated accelerators, and high-end AVX-512 processors.

%%%%%%%%%%%%%%%%%%%%%%%%%%
\section{Work Organization}

This work is organized into five chapters. Chapter 1 presents the context of the project, along with its motivation and objectives. Chapter 2 introduces the CKKS scheme and describes its main operations. Chapter 3 presents the methodology adopted throughout the project, while Chapter 4 defines its functional and non-functional requirements. Finally, Chapter 5 describes the technologies employed, presents the implementation and experimental results, and provides a comparison with the state of the art.

